% Generated by scripts/build_reports.py; edit reviewed Markdown inputs.
% Compile with XeLaTeX.
% Options for packages loaded elsewhere
\PassOptionsToPackage{unicode}{hyperref}
\PassOptionsToPackage{hyphens}{url}
\PassOptionsToPackage{dvipsnames,svgnames,x11names}{xcolor}
\documentclass[
  11pt,
]{article}
\usepackage{xcolor}
\usepackage[margin=1in]{geometry}
\usepackage{amsmath,amssymb}
\setcounter{secnumdepth}{-\maxdimen} % remove section numbering
\usepackage{iftex}
\ifPDFTeX
  \usepackage[T1]{fontenc}
  \usepackage[utf8]{inputenc}
  \usepackage{textcomp} % provide euro and other symbols
\else % if luatex or xetex
  \usepackage{unicode-math} % this also loads fontspec
  \defaultfontfeatures{Scale=MatchLowercase}
  \defaultfontfeatures[\rmfamily]{Ligatures=TeX,Scale=1}
\fi
\usepackage{lmodern}
\ifPDFTeX\else
  % xetex/luatex font selection
\fi
% Use upquote if available, for straight quotes in verbatim environments
\IfFileExists{upquote.sty}{\usepackage{upquote}}{}
\IfFileExists{microtype.sty}{% use microtype if available
  \usepackage[]{microtype}
  \UseMicrotypeSet[protrusion]{basicmath} % disable protrusion for tt fonts
}{}
\makeatletter
\@ifundefined{KOMAClassName}{% if non-KOMA class
  \IfFileExists{parskip.sty}{%
    \usepackage{parskip}
  }{% else
    \setlength{\parindent}{0pt}
    \setlength{\parskip}{6pt plus 2pt minus 1pt}}
}{% if KOMA class
  \KOMAoptions{parskip=half}}
\makeatother
\usepackage{longtable,booktabs,array}
\usepackage{caption}
\captionsetup[table]{skip=6pt}
\newcounter{none} % for unnumbered tables
\usepackage{calc} % for calculating minipage widths
% Correct order of tables after \paragraph or \subparagraph
\usepackage{etoolbox}
\makeatletter
\patchcmd\longtable{\par}{\if@noskipsec\mbox{}\fi\par}{}{}
\makeatother
% Allow footnotes in longtable head/foot
\IfFileExists{footnotehyper.sty}{\usepackage{footnotehyper}}{\usepackage{footnote}}
\makesavenoteenv{longtable}
\setlength{\emergencystretch}{3em} % prevent overfull lines
\providecommand{\tightlist}{%
  \setlength{\itemsep}{0pt}\setlength{\parskip}{0pt}}
\usepackage{amsmath,amssymb}
\usepackage{xurl}
\usepackage{fancyhdr}
\pagestyle{fancy}
\fancyhf{}
\fancyhead[L]{\small Pure Thought | Revised research specifications}
\fancyfoot[L]{\small 2026-09-17}
\fancyfoot[R]{\small\thepage}
\setlength{\headheight}{15pt}
\setlength{\emergencystretch}{3em}
\allowdisplaybreaks
\usepackage{bookmark}
\IfFileExists{xurl.sty}{\usepackage{xurl}}{} % add URL line breaks if available
\urlstyle{same}
% fallback for those not using the hyperref driver hyperxmp:
\makeatletter
\@ifundefined{xmpquote}{\newcommand{\xmpquote}[1]{#1}}{}
\makeatother
\hypersetup{
  pdftitle={10. Flatness--locality--geometry tradeoffs for Chern bands},
  pdfauthor={Pure Thought research portfolio},
  colorlinks=true,
  linkcolor={blue},
  filecolor={Maroon},
  citecolor={Blue},
  urlcolor={blue},
  pdfcreator={LaTeX via pandoc}}

\title{10. Flatness--locality--geometry tradeoffs for Chern bands}
\author{Pure Thought research portfolio}
\date{Revised 17 September 2026}

\begin{document}
\maketitle
\thispagestyle{empty}
\clearpage

\section{10. Flatness--locality--geometry tradeoffs for Chern bands}

\textbf{Revised specification --- 2026-09-11.}
\href{https://github.com/rick-stevens-ai/Pure-Thought/blob/main/reviews/10-Flat-Chern-Bands.md}{Detailed
critique} \ensuremath{\cdot}
\href{https://github.com/rick-stevens-ai/Pure-Thought/blob/main/archive/original-PRDs/10-Flat-Chern-Bands.md}{Original
description} \ensuremath{\cdot}
\href{https://github.com/rick-stevens-ai/Pure-Thought/blob/main/reviews/EVIDENCE-STANDARD.md}{Evidence
standard}

\subsection{Plain-language guide}

\subsubsection{The problem in everyday terms}

In an ordinary energy band, particles can have quite different energies
depending on how they move. A nearly flat band reduces those energy
differences, potentially making interactions between particles more
influential. We want that band to also have a nontrivial topological
structure. The problem asks how well these goals can be combined when
particles are allowed to hop only a limited distance and the model has a
limited number of orbitals.

\subsubsection{Key terms}

\begin{itemize}
\item
  \textbf{Flat band} --- A band whose energy is constant across momentum
  space; a nearly flat band has a small but nonzero variation.
\item
  \textbf{Bandwidth W} --- The difference between the largest and
  smallest energy in the selected band.
\item
  \textbf{Gap \ensuremath{\Delta}} --- The minimum energy separation
  between that band and other bands.
\item
  \textbf{Flatness ratio W/\ensuremath{\Delta}} --- A way to compare
  bandwidth to isolation from other bands; smaller values mean a
  flatter, well-separated band.
\item
  \textbf{Hopping range} --- The greatest lattice distance over which
  the model allows direct hopping.
\item
  \textbf{Chern band} --- A band with a nonzero Chern number, indicating
  nontrivial global topology.
\item
  \textbf{Berry curvature} --- A local measure of how quantum states
  twist as momentum changes; its integral determines the Chern number in
  the stated convention.
\item
  \textbf{Quantum metric} --- A measure of how distinguishable nearby
  momentum-dependent quantum states are.
\item
  \textbf{Spectral flattening} --- A mathematical change to band
  energies that preserves the chosen states but generally changes the
  spatial range of hopping.
\end{itemize}

\subsubsection{Why this matters}

Flat topological bands are useful settings for studying collective
quantum behavior driven by interactions. But there are tradeoffs:
restricting the model's spatial range changes which flatness properties
are possible. Quantifying those tradeoffs can show whether a proposed
improvement is attainable or requires a more complicated model.

\subsubsection{What progress would mean}

Progress would provide certified examples and limits within a defined
design family. Favorable single-particle geometry would make a model
interesting to investigate further; it would not by itself prove an
interacting fractional topological phase.

\subsection{Core question}

For fixed orbital count, hopping range, and gap normalization, how small
can a Chern band's bandwidth and geometric nonuniformity be, with
certified bounds?

\subsection{Scope and assumptions}

Choose a two-dimensional isolated rank-one band, target C=1, and a
compact coefficient family with range R and fixed orbital positions.
Initially use a specified two-band family extending Challenge 09. Impose
a positive gap floor and a norm bound to prevent trivial energy
rescaling.

Maintain two distinct tracks: finite-range approximate flatness, and
exact spectral flattening with generally non-finite-range hopping. Do
not require exact nonzero-Chern flatness with strictly finite-range
hopping under the no-go theorem's hypotheses.

\subsection{Mathematical target}

Define W=max\_k E\_n\ensuremath{-}min\_k E\_n and
\ensuremath{\Delta}=min\_k
min\_\{j\ensuremath{\ne}n\}|E\_j\ensuremath{-}E\_n|. Optimize
W/\ensuremath{\Delta}. For the fixed Brillouin-zone coordinates define
g\_ij=Re\ensuremath{\langle}\ensuremath{\partial}\_iu|(1\ensuremath{-}P)|\ensuremath{\partial}\_ju\ensuremath{\rangle}
and F with the same Chern convention as Challenge 09. Then
C=(1/2\ensuremath{\pi})\ensuremath{\int}F and
\ensuremath{\bar{F}}=2\ensuremath{\pi}C/Area(BZ).

Report separately
Var(F)=Area\ensuremath{{}^{-1}}\ensuremath{\int}(F\ensuremath{-}\ensuremath{\bar{F}})\ensuremath{{}^{2}}
and D\_g=Area\ensuremath{{}^{-1}}\ensuremath{\int}(tr
g\ensuremath{-}|F|). These are diagnostics under a specified Euclidean
coordinate metric, not a universal fractional-Chern-insulator stability
score.

\subsection{Required outputs and proof obligations}

Export coefficient/domain constraints, W upper and \ensuremath{\Delta}
lower bounds, certified C, geometric enclosures, and optimization gaps.
An achieved objective gives an upper bound on the minimum; a global
lower bound needs a covering/relaxation proof. For infinite-range
flattening, give decay and truncation estimates.

\subsection{Validation and rejection controls}

Verify normalization \ensuremath{\int}F=2\ensuremath{\pi}C and the
pointwise metric inequality in the chosen convention. Check that
truncating a spectrally flattened model generally reintroduces
bandwidth. Include a trivial exactly flat band as a negative topology
test. Compare candidate objectives using identical units, embedding, and
norm constraints.

\subsection{Milestones and research extension}

Baseline: one certified near-flat Chern model and a correct locality
audit. Strong: global objective bounds on one compact family or a
rigorous truncation tradeoff. Research extension: optimal geometry
within a specified class; many-body FCI stability requires an added
interacting Hamiltonian and filling.

Milestones are acceptance gates, not calendar promises. Before a run,
freeze the concrete instance/domain, tolerance, compute budget, and
permitted dependencies. A bounded search that fails to find a result
must preserve an unresolved outcome.

\subsection{Starting instruction}

Work on the baseline above. First make every input and convention
explicit, then validate the reference case. Build the checker alongside
the calculation. Label numerical evidence, finite-model certificates,
and continuum theorems separately; return the strongest justified result
and identify the exact unmet proof obligation.

\subsection{Primary references and their role}

\begin{itemize}
\tightlist
\item
  \href{https://arxiv.org/abs/1311.4956}{Chen et al., Flat Chern-band
  locality obstruction} --- No-go theorem for simultaneous exact
  flatness, nonzero Chern number, and strictly local hopping.
\end{itemize}

References support definitions or limitations; the scope choices and
proposed milestones are editorial recommendations, not claims that these
results have already been obtained.

\subsection{Displayed formulation and claim boundary}

\href{https://github.com/rick-stevens-ai/Pure-Thought/blob/main/PRDs/10-Flat-Chern-Bands.md}{Canonical
specification}.

\[W=\max_k E_n(k)-\min_k E_n(k),\quad
\Delta=\min_k\min_{j\ne n}|E_j(k)-E_n(k)|,\quad\text{minimize }W/\Delta.\]

\[g_{ij}=\operatorname{Re}\langle\partial_i u|(1-P)|\partial_j u\rangle,
\quad \bar F=\frac{2\pi C}{A},\quad A=\operatorname{Area}(\mathrm{BZ}),\]

\[\operatorname{Var}(F)=\frac1A\int_{\mathrm{BZ}}(F-\bar F)^2d^2k,
\qquad D_g=\frac1A\int_{\mathrm{BZ}}(\operatorname{tr}g-|F|)d^2k.\]

Fix a gap floor, norm bound, orbital positions and coordinate metric.
Use the Chern convention of 09. Strictly finite-range approximate
flatness and generally long-range exact flattening are separate tasks;
these diagnostics alone do not prove an interacting fractional phase.

\subsection{Detailed review and rationale}

\href{https://github.com/rick-stevens-ai/Pure-Thought/blob/main/PRDs/10-Flat-Chern-Bands.md}{Revised
problem} \ensuremath{\cdot}
\href{https://github.com/rick-stevens-ai/Pure-Thought/blob/main/archive/original-PRDs/10-Flat-Chern-Bands.md}{Preserved
original} \ensuremath{\cdot}
\href{https://github.com/rick-stevens-ai/Pure-Thought/blob/f854d4e12b55c1f5b5c2a964b75f2d23496d981e/PRDs/10-Flat-Chern-Bands.md}{Pinned
upstream source}

Line references below refer to the preserved original at commit
\texttt{f854d4e12b55}, not the rewritten file. ``Blocking'' means the
original issue can invalidate the scientific target or claimed
certificate; ``major'' means it materially changes scope, inference, or
acceptance criteria. These are critiques of specifications, not results
of running the embedded research-code sketches.

\subsubsection{Detailed findings}

\begin{enumerate}
\def\labelenumi{\arabic{enumi}.}
\item
  \textbf{Blocking --- incompatible design requirements (33--40;
  synthesis finite-range search).} Exact flatness, nonzero Chern number,
  and strictly local finite-range hopping cannot all be required under
  the cited theorem's assumptions. The rewrite splits the problem into
  two valid tracks.
\item
  \textbf{Blocking --- normalization (38, 86).} The two displayed
  constant-curvature formulas disagree and omit/misplace
  2\ensuremath{\pi}. Defining C by its integral fixes
  \ensuremath{\bar{F}} uniquely.
\item
  \textbf{Major --- unsupported objective (22--24, 39--40).} The stated
  stability ratio is not a general criterion proving an FCI ground
  state. Geometry, bandwidth, and an interacting many-body gap are
  different claims.
\item
  \textbf{Major --- exact versus approximate (787--805).} A tiny
  numerical bandwidth or curvature variance does not establish exact
  flatness or uniformity. The revision asks for certified enclosures and
  distinguishes achievable values from global optima.
\end{enumerate}

\subsubsection{What the revision preserves and changes}

The central motivation is retained, but the executable question is now:
\textbf{For fixed orbital count, hopping range, and gap normalization,
how small can a Chern band's bandwidth and geometric nonuniformity be,
with certified bounds?} The revised scope chooses a concrete starting
model; those choices are proposed research-design decisions, not facts
inferred from the original. The baseline can succeed without resolving
the full open research program.

The replacement also removes the original implementation sketches from
the active specification. They remain in the archive for comparison, but
have not been made executable or certified. New implementations should
derive their equations from the revised definitions rather than copy
unchecked prototypes.

\subsubsection{Acceptance criteria to use instead}

Baseline: one certified near-flat Chern model and a correct locality
audit. Strong: global objective bounds on one compact family or a
rigorous truncation tradeoff. Research extension: optimal geometry
within a specified class; many-body FCI stability requires an added
interacting Hamiltonian and filling.

\subsubsection{Sources supporting the corrections}

\begin{itemize}
\tightlist
\item
  \href{https://arxiv.org/abs/1311.4956}{Chen et al., Flat Chern-band
  locality obstruction} --- No-go theorem for simultaneous exact
  flatness, nonzero Chern number, and strictly local hopping.
\end{itemize}

\section{Evidence and execution standard}

This standard applies to the revised specifications. It is a proposed
contract for future research runs; this editorial revision does not
supply the scientific certificates themselves.

\subsection{What ``pure thought'' means here}

Use explicitly stated mathematical models and deductions without fitting
conclusions to experimental outcomes. Synthetic examples, exact model
inputs, analytic references, numerical search, and established software
are allowed. Physical constants, electronic-structure approximations,
constitutive laws, and empirical parameter sets must be identified when
used. A model-derived result is conditional on that model.

The original ban on legacy software conflicts with its proposed reliance
on NumPy, SciPy, CVXPY, nauty, and other libraries. Reimplementing every
optimizer and algebra system is not necessary for scientific
independence. Prefer a small independently checkable certificate and a
documented trusted base. A fresh implementation can be a project
objective, but must not be confused with a mathematical proof or novelty
claim.

\subsection{Four evidence levels}

{\def\LTcaptype{none} % do not increment counter
\begin{longtable}[]{@{}
  >{\raggedright\arraybackslash}p{(\linewidth - 4\tabcolsep) * \real{0.3333}}
  >{\raggedright\arraybackslash}p{(\linewidth - 4\tabcolsep) * \real{0.3333}}
  >{\raggedright\arraybackslash}p{(\linewidth - 4\tabcolsep) * \real{0.3333}}@{}}
\toprule\noalign{}
\begin{minipage}[b]{\linewidth}\raggedright
Label
\end{minipage} & \begin{minipage}[b]{\linewidth}\raggedright
What it establishes
\end{minipage} & \begin{minipage}[b]{\linewidth}\raggedright
What it does not establish
\end{minipage} \\
\midrule\noalign{}
\endhead
\bottomrule\noalign{}
\endlastfoot
Numerical evidence & A reproducible computation at stated precision,
sample size, and discretization & Exact truth, continuum coverage, or a
universal theorem \\
Finite-model certificate & An exact/validated statement about the
specified finite model or finite relaxation & An omitted tail,
thermodynamic limit, or physical realization \\
Conditional mathematical theorem & A proof over the stated domain with
every hypothesis and error controlled & Truth outside those assumptions
or a demonstrated physical system \\
Formal proof & A checked theorem in a named proof assistant, with
imported axioms/dependencies listed & Correct physical modeling merely
by virtue of formalization \\
\end{longtable}
}

Formalization is an optional extra deliverable unless the specific
project makes it mandatory. High precision, JSON output, solver status
and small residuals do not by themselves confer any of these proof
levels.

\subsection{Claim directions}

For minimization over a larger feasible set, the optimum is a lower
bound on the original minimum. A physical feasible point supplies an
upper bound. For maximization the directions reverse. Track feasibility
of the primal/dual witnesses separately from their objective values.

Infeasibility of an inner ansatz excludes only that ansatz.
Infeasibility of an outer relaxation can exclude the original problem if
the map into the relaxation is proved and all approximations preserve
necessity. Passing necessary tests means ``not excluded''; existence
needs a constructive sufficient argument.

A smallest observed example is not a minimality theorem. A complete
catalog must define the universe, equivalence relation, boundaries, and
proof that pruning misses nothing. A finite-size numerical trend does
not prove an asymptotic limit.

\subsection{Minimum run manifest}

Before calculation, record the problem ID and revision hash, exact
inputs and conventions, quantified domain, equivalence relation, target
observable, accepted evidence level, tolerances, and resource limits.
Record dependencies and precision/rounding mode. Where randomization is
used, preserve seeds, trial counts and raw sufficient statistics.
Resource limits are planning decisions, not evidence that a mathematical
solution cannot exist.

An output bundle should include:

\begin{itemize}
\tightlist
\item
  The statement being claimed, including all assumptions and
  quantifiers.
\item
  Construction or proof object, plus an independently runnable verifier
  and its trust assumptions.
\item
  Numerical/error ledger separating roundoff, discretization,
  truncation, optimization, sampling and modeling error where relevant.
\item
  Positive reference cases and negative controls that the verifier must
  reject.
\item
  Coverage: included domain, omitted domain, unresolved regions and
  termination reason.
\end{itemize}

These are requirements for future implementations, not a promise of
polynomial-size certificates for every problem.

\subsection{Checking specific certificate types}

\textbf{LP/SDP/SOS:} verify equality constraints, objective direction,
and cone feasibility using exact arithmetic or outward-rounded
intervals. Rationalizing a numerical solution can break PSD constraints.
Correct residuals with justified norm bounds; do not silently project
and reuse the old objective. Distinguish a proof of a polynomial
approximation from a proof for the original function.

\textbf{Interval numerics:} state function enclosures, domains and
outward rounding. A small root box needs existence and, if claimed,
uniqueness. Cover the complement for a global enumeration. Include
singularities and tails rather than discarding them numerically.

\textbf{SAT/SMT:} a solver's UNSAT label is not a certificate. Preserve
a proof supported by a checker and the proof that the encoding, symmetry
reduction and blocked solutions match the mathematical universe. DRAT
applies to an appropriate Boolean CNF proof, not automatically to every
real-arithmetic or SDP claim.

\textbf{Statistics:} specify the estimand, dependence assumptions,
confidence procedure and bias. Zero observed failures does not prove
zero failure probability. MCMC error bars based on estimated
autocorrelation are not generally rigorous mixing guarantees.

\subsection{Milestones and novelty}

Baseline means a working, independently validated result on a concrete
reference problem. Strong means a broader or sharper result with its
claimed evidence level. Research extension means a possible novel result
whose feasibility and significance need review. A negative theorem can
be valuable; an unsuccessful search can also be useful if its scope is
documented, but the two are not interchangeable.

Do not promise ``publication-quality'' in a fixed number of months.
Before claiming novelty, conduct a focused current literature comparison
against the exact theorem/model/metric. The sources in this revision
check key definitions and limitations; they are not an exhaustive
novelty audit across all 30 fields.

\end{document}
